\addtocontents{toc}{\protect\setcounter{tocdepth}{0}}
\subsection{Spektralsatz der Linearen Algebra} \label{sec:spektralsatzlinearealgebra}
\addtocontents{toc}{\protect\setcounter{tocdepth}{3}}

Da der Spektralsatz eine zentrale Aussage in der Linearen Algebra ist, verweisen wir für den Beweis auf die zahlreiche Literatur, wie zum Beispiel \cite[Kap.~6.7]{fischer_lineare_algebra_20}. 

\begin{theorem} \label{thm:spektralsatzderlinearenalgebra}
Jede reelle, symmetrische $n \times n$-Matrix $A$ ist orthogonal diagonalisierbar.
\end{theorem}
Das bedeutet:
\begin{itemize}
    \item Alle Eigenwerte von $A$ sind reell.
    \item Es existiert eine Orthonormalbasis des $\mathbb{R}^n$, die ausschließlich aus Eigenvektoren von $A$ besteht.
    \item Die Matrix lässt sich zerlegen in $A = Q D Q^T$, wobei $D$ eine Diagonalmatrix der Eigenwerte und $Q$ eine orthogonale Matrix der Eigenvektoren ist.
\end{itemize}

In unserem Fall haben wir die Matrix 
\begin{align}
    M := \sum_{w \in \Omega} w(\vec{s}) w(\vec{s})^T \in \mathbb{R}^{3 \times 3}, \nonumber 
\end{align}
welche symmetrisch ist. Dies lässt sich unter Ausnutzung der Rechenregeln für transponierte Matrizen ($(A+B)^T = A^T + B^T$ und $(AB)^T = B^T A^T$) direkt zeigen:
\begin{align}
    M^T &= \left( \sum_{w \in \Omega} w(\vec{s}) w(\vec{s})^T \right)^T \nonumber \\
    &= \sum_{w \in \Omega} \left( w(\vec{s}) w(\vec{s})^T \right)^T \nonumber \\
    &= \sum_{w \in \Omega} (w(\vec{s})^T)^T w(\vec{s})^T \nonumber \\
    &= \sum_{w \in \Omega} w(\vec{s}) w(\vec{s})^T = M. \nonumber
\end{align}